\documentclass[11pt]{article}

\usepackage{geometry}
\geometry{a4paper, margin=2.4cm}

\usepackage{fontspec}
\usepackage{amsmath}
\usepackage{amsthm}
\usepackage{unicode-math}
% Fonts: Latin Modern (fontspec / unicode-math defaults) — bundled with tectonic.

\usepackage{microtype}
\usepackage[dvipsnames]{xcolor}
\usepackage[colorlinks=true, linkcolor=RoyalBlue, urlcolor=RoyalBlue]{hyperref}
\usepackage{enumitem}

\setlength{\parindent}{0pt}
\setlength{\parskip}{0.45em}
\setlength{\emergencystretch}{3em}
\setlist{leftmargin=1.4em}

\theoremstyle{plain}
\newtheorem{theorem}{Theorem}
\newtheorem{proposition}{Proposition}

\newcommand{\tp}{^{\top}}
\newcommand{\bc}{\mathbf{c}}
\newcommand{\bb}{\mathbf{b}}
\newcommand{\bG}{\mathbf{G}}
\newcommand{\bZ}{\mathbf{Z}}
\newcommand{\bA}{\mathbf{A}}
\newcommand{\bd}{\mathbf{d}}
\newcommand{\bu}{\mathbf{u}}
\newcommand{\by}{\mathbf{y}}
\newcommand{\bs}{\mathbf{s}}
\newcommand{\bw}{\mathbf{w}}
\newcommand{\bze}{\mathbf{z}}
\newcommand{\bone}{\mathbf{1}}

\title{\bfseries On the support size of the genomic OCS optimum}
\author{\textbf{Adel Kaleche}\\[2pt] \small\itshape Affiliation to be completed.}
\date{}

\begin{document}
\maketitle

\begin{abstract}
\noindent
A companion note to the support-first solver. For genomic optimum contribution selection
(OCS) we bracket the size of the optimal support: a clean $n$-independent bound holds in the
no-ridge limit (Theorem~\ref{thm:eps0}), no universal bound survives the ridge
(Theorem~\ref{thm:noridge}), and the practical regime between them is governed by the joint
geometry of spectrum, breeding values and cap, admitting no single-scalar law. A KKT identity
(Proposition~\ref{prop:kkt}) connects the two theorems. All empirical claims are reproducible.
\end{abstract}

\section{The question}

Genomic OCS solves, for $n$ candidates with breeding values $\bb \in \mathbb{R}^n$ and
relationship matrix $\bG = \bZ\bZ\tp/s + \varepsilon\mathbf{I}$ ($\bZ$ the $n\times m$ centred
genotype matrix, $s$ a scale, $\varepsilon$ a small ridge for positive-definiteness),
\begin{equation}\label{eq:ocs}
  \max\ \bb\tp\bc \quad\text{s.t.}\quad \bA\bc = \bd,\ \ \mathbf{0} \le \bc \le \bu,\ \ \bc\tp\bG\bc \le k,
\end{equation}
where $\bA$ is the $q\times n$ budget matrix ($q=1$ for the simplex $\bone\tp\bc=1$; $q=2$ for
the sexed form $\sum_{\text{males}} = \sum_{\text{females}} = \tfrac12$). Empirically the
optimum $\bc^\star$ activates very few candidates: its support $S = \{i : c^\star_i > 0\}$ is
$\approx 15$--$30$ on real panels and stays bounded as $n$ grows at a fixed cap. What bounds
$|S|$?

The answer is a \emph{bracket}. In the no-ridge limit there is a clean $n$-independent bound
(\S\ref{sec:eps0}); once the ridge is present there is provably no universal bound, the
support reaching the whole population (\S\ref{sec:noridge}); and the practical regime between
them is set by the joint geometry of $(\text{spectrum}, \bb, k)$, which admits no
single-scalar law and which we map empirically (\S\ref{sec:empirical}). A KKT identity
(\S\ref{sec:kkt}) bridges the two theorems.

\section{The \texorpdfstring{$\varepsilon = 0$}{eps=0} bound}\label{sec:eps0}

\begin{theorem}\label{thm:eps0}
For $\varepsilon = 0$ (so $\bG = \bG_0 = \bZ\bZ\tp/s$ of rank $r \le m$),
\eqref{eq:ocs} has an optimal $\bc^\star$ with $|S| \le q + r + 1$, independent of $n$.
\end{theorem}

\begin{proof}
By strong duality (Slater holds at any interior-feasible instance) there is a multiplier
$\lambda^\star \ge 0$ such that $\bc^\star$ maximises the Lagrangian
$L(\bc) = \bb\tp\bc - \lambda^\star\,\bc\tp\bG_0\bc$ over the polytope
$\{\bA\bc = \bd,\ \bc \ge 0\}$. Since $\bc\tp\bG_0\bc = \lVert\bZ\tp\bc\rVert^2/s$, the value
$L(\bc)$ depends on $\bc$ only through the pair $(\bZ\tp\bc,\ \bb\tp\bc) \in \mathbb{R}^{r+1}$.
Hence the slice
\[
  P = \{\, \bc \ge 0 : \bA\bc = \bd,\ \bZ\tp\bc = \bZ\tp\bc^\star,\ \bb\tp\bc = \bb\tp\bc^\star \,\}
\]
is \emph{entirely optimal}: every $\bc \in P$ is feasible ($\bc\tp\bG_0\bc = \lVert\bZ\tp\bc^\star\rVert^2/s \le k$)
and attains the optimal value $\bb\tp\bc^\star$. $P$ is nonempty (it contains $\bc^\star$) and
bounded (the budget together with $\bc \ge 0$), and it is cut by $q + r + 1$ equality rows, so
it has a vertex; a vertex of $\{\bc \ge 0 : \mathbf{M}\bc = \mathbf{e}\}$ with $q+r+1$ rows has
at most $q + r + 1$ nonzeros. That vertex is an optimum with the stated support.
\end{proof}

This is the Carathéodory / Barvinok--Pataki idea applied correctly: one does \emph{not} count
on the curved boundary of the ellipsoid (where, $\bG$ being positive definite, every point is
already an extreme point of the feasible set, so the support is unconstrained) --- instead one
\emph{fixes the low-rank image $\bZ\tp\bc$ and the objective}, which makes the optimal slice
affine, and counts LP vertices there. Active-set solvers (support-first; the critical line
algorithm) return such a vertex, whereas interior-point and ADMM solvers return non-sparse
interior points and threshold post hoc --- so the bound describes precisely what an active-set
solver computes.

\section{No universal bound for \texorpdfstring{$\varepsilon > 0$}{eps>0}}\label{sec:noridge}

\begin{theorem}\label{thm:noridge}
For $\varepsilon > 0$ no bound on $|S|$ of the form $f(q, r)$ independent of $n$ exists: the
support can equal $n$.
\end{theorem}

\begin{proof}[Proof (counterexample)]
Take $\bG = \varepsilon\mathbf{I}$ --- the degenerate case $m = 0$ (no markers, every pair
equally unrelated; $\operatorname{rank}(\bG_0) = 0$). Then \eqref{eq:ocs}, simplex form, is
\[
  \max\ \bb\tp\bc \quad\text{s.t.}\quad \bone\tp\bc = 1,\ \bc \ge 0,\ \varepsilon\lVert\bc\rVert^2 \le k,
\]
and $\lVert\bc\rVert^2 = \sum_i c_i^2$ is exactly the rate-of-inbreeding proxy. The minimum of
$\lVert\bc\rVert^2$ on the simplex is $1/n$, attained only at the uniform plan $\bc = \bone/n$.
For $k$ slightly above $\varepsilon/n$ the feasible set is the simplex intersected with a ball
that shrinks onto $\bone/n$, forcing $|S| = n$. Numerically (\texttt{bound\_validation.py},
block~7, $n = 300$): $|S| = 300, 295, 256, 151, 52, 12$ as $k$ loosens from $1.05\times$ to
$50\times$ the minimum. So $|S|$ ranges over $[1, n]$ with no $n$-independent ceiling.
\end{proof}

Genetically, $\bG = \varepsilon\mathbf{I}$ is the \emph{no-structure} limit; with no relatedness
pattern the diversity cap can be met only by spreading contributions over the whole population.
The small support seen on real data is therefore a property of the \emph{structure} of $\bG$
(and of $\bb$), not a worst-case guarantee.

\section{The bridge: a KKT identity}\label{sec:kkt}

\begin{proposition}\label{prop:kkt}
At an optimum with the kinship constraint active, the contributions on the support are an
affine function of each candidate's augmented feature $(b_i, \bze_i) \in \mathbb{R}^{m+1}$
($\bze_i$ the $i$-th row of $\bZ$): $\ c^\star_i = \alpha\,b_i + \bw\tp\bze_i + \beta_{\mathrm{sex}(i)}$
for $i \in S$.
\end{proposition}

\begin{proof}
KKT gives $\mu \in \mathbb{R}^q$, $\lambda \ge 0$, $\bs \ge 0$ with
$\bb = \bA\tp\mu + 2\lambda\bG\bc^\star - \bs$ and $s_i c^\star_i = 0$. On $S$ ($s_i = 0$):
$b_i = (\bA\tp\mu)_i + 2\lambda(\bG\bc^\star)_i$. Writing $\by := \bZ\tp\bc^\star$ and using
$\bG\bc^\star = \bZ\by/s + \varepsilon\bc^\star$ gives, for $i \in S$,
$\varepsilon c^\star_i = (b_i - (\bA\tp\mu)_i)/(2\lambda) - (\bze_i\tp\by)/s$, i.e.
$c^\star_i = \alpha\,b_i + \bw\tp\bze_i + \beta_{\mathrm{sex}(i)}$ with $\alpha = 1/(2\lambda\varepsilon)$,
$\bw = -\by/(s\varepsilon)$, $\beta_{\mathrm{sex}(i)} = -(\bA\tp\mu)_i/(2\lambda\varepsilon)$.
\end{proof}

The support's contributions thus live on an $(m+2)$-parameter family. As $\varepsilon \to 0$
the $\varepsilon c^\star_i$ term vanishes and the identity forces $\bb - \bA\tp\mu$ into the
rank-$r$ row space of $\bZ$, recovering Theorem~\ref{thm:eps0} and exhibiting $m$ (the marker
count) as the relevant dimension. For $\varepsilon > 0$ the same identity does \emph{not} cap
$|S|$, consistent with Theorem~\ref{thm:noridge}.

\section{The empirical regime between the theorems}\label{sec:empirical}

Between the two clean statements lies the regime breeders use, mapped across synthetic spectra,
$\bb$-alignments and caps and on the real panels:

\begin{itemize}
\item \textbf{$\bb$'s alignment with the spectrum drives $|S|$, inversely.} With $\bG$ and $k$
  fixed, putting $\bb$ on the dominant (coancestry-expensive) eigendirection forces a large
  support (mean $|S| \approx 146$ of $n = 800$ for $\bb$ the top eigenvector); spreading $\bb$
  across many cheap directions collapses it to $\approx 4$. Gain sought in an expensive
  direction must be diluted over many candidates to respect the cap (\texttt{bound\_balign.py}).
\item \textbf{No single scalar predicts $|S|$ across regimes.} Neither $\operatorname{rank}(\bG_0)$,
  the effective rank / participation ratio of the spectrum, the directional cost
  $\bb\tp\bG\bb/\bb\tp\bb$, nor the support of the unconstrained direction $(\bG^{-1}\bb)_+$
  tracks $|S|$ across both synthetic and real instances; the best dimensionless combination,
  $|S| \sim (\bb\tp\bG\bb/\bb\tp\bb \cdot 1/k)^{0.8}$, reaches only $R^2 \approx 0.58$ and
  mispredicts the real panels' relative support (\texttt{bound\_predictor.py},
  \texttt{bound\_lawfit.py}).
\item \textbf{The curve $|S|(k)$ is a per-instance power law of non-universal exponent.}
  $|S| \sim k^{-\alpha}$ fits each instance well ($R^2 \approx 0.9$--$0.98$ where the support is
  well resolved), with $\alpha \approx 1$ on both real panels but ranging $0.5$--$1.7$ across
  synthetic spectra and alignments (\texttt{bound\_curve.py}).
\item \textbf{Real panels.} Wheat ($n = 599$) and mouse ($n = 1814$): $|S| = 25/26$ and $17/26$
  (simplex / sexed-optiSel) at the working cap, both of order tens --- small and stable in $n$,
  as the decaying-spectrum regime predicts qualitatively, though no formula pins the value
  (\texttt{bound\_real.py}).
\end{itemize}

Theorem~\ref{thm:noridge} explains \emph{why} no scalar law exists: there is no universal bound
to predict, so the practical value is a function of the full $(\text{spectrum}, \bb, k)$
geometry.

\section{A conditional bound: the obstruction and the conjecture}

The universal question is closed negatively (Theorem~\ref{thm:noridge}); the useful one is
\emph{conditional}. Two clean bounds bracket the rest: at a \emph{loose cap} ($\bc\tp\bG\bc^\star < k$)
the active set is the polytope $\{\bA\bc=\bd,\ \bc\ge0\}$ alone and $\bc^\star$ is an LP vertex, so
$|S| \le q$; at \emph{exact rank} $d$, Theorem~\ref{thm:eps0} gives $|S| \le q + d + 1$.

A bound in terms of \emph{effective} rank --- the realistic ask, since real spectra decay --- does
not follow from the same argument, and the obstruction is precise. Theorem~\ref{thm:eps0} linearises
the optimal slice by pinning $\bZ\tp\bc$, the \emph{entire} range of $\bG_0$; for a full-rank $\bG$
one would have to pin all $n$ projections $u_i\tp\bc$, giving only the trivial $|S| \le q + n + 1$.
One cannot pin just the top-$d$ projections and drop the tail, because the tail contributes
$\bc\tp E\bc = \sum_{i>d} \lambda_i (u_i\tp\bc)^2$ --- a \emph{genuine quadratic}, not a low-rank term
--- so fixing it is not a linear constraint and the slice is not a polytope. The tail's curvature is
exactly what lets the support spread; effective rank cannot enter a vertex count. A conditional bound
therefore needs a different (spectral-perturbation) argument.

It also needs two assumptions, \emph{both necessary}:
\begin{itemize}
\item \textbf{(A1) No large degenerate cheap plateau.} A block of many near-equal small eigenvalues is
  a cheap subspace the optimum spreads into; $\bG = \varepsilon\mathbf{I}$ is the extreme case, giving
  $|S| = n$ (Theorem~\ref{thm:noridge}). A smoothly decaying spectrum has no such plateau.
\item \textbf{(A2) $\bb$ avoids the dominant directions.} Even with a gap, $\bb$ on the top (expensive)
  eigenvector forces dilution and a large support (mean $|S| \approx 146$ of $800$ in the alignment
  sweep).
\end{itemize}
Both counterexamples are established, so neither assumption can be dropped. Under (A1)+(A2) --- the
regime of real genomic matrices, where $|S| \approx 15$--$30$ --- we conjecture $|S|$ is small and
$n$-independent, of the order of the dominant eigendirections $\bb$ actually excites; the growth half
($|S|$ rising as $k \to 0$) follows the
contributions\,$\leftrightarrow\,\Delta F\,\leftrightarrow\,N_e$ link (Wray \& Thompson 1990;
Woolliams \& Bijma 2000). Proving it --- coupling the spectral gap with $\bb$'s projection onto the
top eigenspace, robust to the tail's curvature --- is the open problem, the meeting point of
optimisation (the perturbed-cone face) and quantitative genetics (effective lineages).

\section*{Numerical evidence (reproducible)}

\texttt{bound\_validation.py} (the $\varepsilon=0$ bound, $n$-independence, ridge and
spectral-decay sweeps, and the Theorem~\ref{thm:noridge} counterexample);
\texttt{bound\_real.py} (wheat and mouse, after \texttt{research/repro/*\_export.R});
\texttt{bound\_balign.py} (the $\bb$-alignment sweep); \texttt{bound\_predictor.py} and
\texttt{bound\_lawfit.py} (the predictor / law search); \texttt{bound\_curve.py} (the
per-instance $|S|(k)$ fit). Full log: \texttt{research/support\_bound\_sketch.md}.

\section*{References}
\sloppy
\begin{enumerate}[leftmargin=2em]
\item \textbf{Pataki, G. (1998).} On the rank of extreme matrices in semidefinite programs and the multiplicity of optimal eigenvalues. \emph{Math.\ Oper.\ Res.} 23(2):339--358. DOI 10.1287/moor.23.2.339.
\item \textbf{Markowitz, H. (1956).} The optimization of a quadratic function subject to linear constraints. \emph{Naval Res.\ Logist.\ Q.} 3(1--2):111--133. DOI 10.1002/nav.3800030110.
\item \textbf{Wray, N.R. \& Thompson, R. (1990).} Prediction of rates of inbreeding in selected populations. \emph{Genet.\ Res.} 55(1):41--54. DOI 10.1017/S0016672300025180.
\item \textbf{Woolliams, J.A. \& Bijma, P. (2000).} Predicting rates of inbreeding in populations undergoing selection. \emph{Genetics} 154(4):1851--1864. DOI 10.1093/genetics/154.4.1851.
\item \textbf{Yamashita, M., Mullin, T.J. \& Safarina, S. (2018).} An efficient second-order cone programming approach for optimal selection in tree breeding. \emph{Optim.\ Lett.} 12(7):1683--1697. DOI 10.1007/s11590-018-1229-y.
\item \textbf{Waldmann, P. (2025).} Genomic optimum contribution selection and mate allocation using JuMP. \emph{Bioinform.\ Adv.} 5(1):vbaf259. DOI 10.1093/bioadv/vbaf259.
\end{enumerate}

\end{document}
